\documentclass{article}

\usepackage[utf8]{inputenc}
\usepackage[T1]{fontenc}
\usepackage{amsmath,amssymb,amsthm}
\usepackage{mathtools}
\usepackage{hyperref}
\usepackage{booktabs}

\makeatletter
\newcommand{\citep}[2][]{%
  \begingroup
  \def\@tmpkey{#2}%
  \ifx\@tmpkey\@empty\else
    (\ifx\@empty#1\@empty\else #1, \fi
     \@nameuse{bibkey@\@tmpkey})%
  \fi
  \endgroup
}
\expandafter\def\csname bibkey@lasser2026gfm\endcsname{Paper~1}
\expandafter\def\csname bibkey@lasser2026poset\endcsname{Paper~2}
\expandafter\def\csname bibkey@lasser2026horizon\endcsname{Paper~3}
\expandafter\def\csname bibkey@lasser2026exo\endcsname{Paper~5}
\expandafter\def\csname bibkey@lasser2026phase\endcsname{Paper~6}
\makeatother

\newtheorem{definition}{Definition}
\newtheorem{proposition}{Proposition}
\newtheorem{lemma}{Lemma}
\newtheorem{theorem}{Theorem}
\newtheorem{remark}{Remark}
\newtheorem{assumption}{Assumption}

\DeclareMathOperator{\vol}{vol}
\newcommand{\volL}{{\vol_{\mathrm{P}}}}
\newcommand{\rext}{r_{\mathrm{ext}}}
\newcommand{\meff}{m_{\mathrm{eff}}}
\newcommand{\mindep}{m_{\mathrm{eff}}^{\mathrm{indep}}}
\newcommand{\Vdisc}{V_{\gamma}}
\newcommand{\Tadv}{T_{\mathrm{adv}}}
\newcommand{\Ddiv}{\Delta_{\mathrm{div}}}
\newcommand{\pshock}{p_{\mathrm{shock}}}
\newcommand{\lossDmax}{\ell_D^{\max}}
\newcommand{\lossdivmax}{\ell_{\mathrm{div}}^{\max}}

\title{Draft v4: Lemma 5c Proof Sketch ---\\
  Minimax Static Tightening of the Anti-Monopolar Region\\
  \large{(v2 after codex Round A; v3 after Round B; v4 after Round C:
  single-shock reduction tightened to $N=1$ a.s., multi-shock
  scope-restricted to cooperative-overlap regime, $\pshock$ reframed
  as Poisson rate $\lambda$, removable singularity at $\kappa=1$
  noted, abstract regression fixed)}}
\author{Paper 10 exploratory draft (companion to lemma\_5\_anti\_monopolar\_robustness.tex)}
\date{}

\begin{document}

\maketitle

\begin{abstract}
Three rounds of codex review have shaped this draft.  Round~A on v1
identified four soundness errors and confirmed the qualitative
minimax lift.  Round~B on v2 confirmed the single-shock form as
citeable but flagged the multi-shock $A_{\mathrm{multi}}$ as
dimensionally inconsistent with the single-shock $\Ddiv$.  Round~C
on v3 confirmed v3's discounted-marginal-loss multi-shock fix in
the balanced/additive special case but flagged the general
non-equal substrate case as requiring sequential counterfactual
derivation.

V4 adopts the scope restriction Round~C suggested.  The single-shock
$\Ddiv$ form is unchanged from v3 (citeable per Round~B, with the
counterfactual $\alpha_s$ definition from v3):
\[
  \Ddiv \;=\; \volL_K \cdot (\lossDmax - \lossdivmax) \cdot
  \mathbb{E}[\gamma^{\Tadv}].
\]
The multi-shock $\Ddiv^{\mathrm{multi}}$ is the discounted sum of
marginal loss differences at shock times
(Definition~\ref{def:ddiv_multi}):
\[
  \Ddiv^{\mathrm{multi}} \;=\; \mathbb{E}\!\left[\sum_{i=1}^{\infty}
  \gamma^{T_i} \cdot
  \bigl(\delta_i^D - \delta_i^{\mathrm{div}}\bigr)\right].
\]

Scope: v4 explicitly restricts the multi-shock claim to the
\textbf{cooperative-overlap regime}, where the substrate-cooperative
structure is dominated by cross-substrate cooperatives (Paper~3's
$\rext$ source) rather than redundant capabilities.  In this regime,
the equal-substrate analysis is a \emph{conservative} lower bound on
the actual sequential-counterfactual $\Ddiv^{\mathrm{multi}}$, which
is what deployment-safety claims need.  In the redundancy regime
(capabilities replicated across substrates), the equal-substrate
analysis is anti-conservative; sequential counterfactual derivation
is open work and Lemma~5c does not cover it.

The canonical tripartite identification (Human + AI +
Formal-Operational) sits in the cooperative-overlap regime by
design: cross-substrate cooperatives (verified workflow, governance
enforcement) are the source of the deployment's $\rext$, not
redundant capabilities.  Paper~10 must include a
\emph{cooperative-vs-redundancy audit} in deployment tooling to
verify the regime; the canonical identification passes by
construction.

Other v4 cleanups:
\begin{itemize}
\item Single-shock reduction tightened from ``$T_i = \infty$ for
  $i \geq 2$ with negligible probability'' to formal $N = 1$ a.s.
  (Round~C Q1).
\item $\pshock$ reframed as continuous-time Poisson rate $\lambda$
  rather than per-period probability (Round~C Q3); $\kappa$ formula
  changes if a discrete-time model is used.
\item The closed form for $\Ddiv^{\mathrm{multi}}$ at $\kappa = 1$
  has a removable singularity (Round~C Q3); the regime condition for
  non-trivial safe region is $\kappa < 1 - \delta$ for some $\delta$,
  not $\pshock$ bounded.
\item Abstract regression fixed (v3 abstract still referenced
  $A_{\mathrm{multi}}$ survival-stream; v4 abstract reflects
  marginal-loss formulation).
\end{itemize}

Status: v4's single-shock form is citeable (per Round~B). v4's
multi-shock form is citeable in the cooperative-overlap regime (per
Round~C with v4's scope restriction). Sequential counterfactual
derivation for the redundancy regime is named as open work.
\end{abstract}

\section{What Lemma 5c claims}

\begin{lemma}[Lemma 5c, target --- v2]
\label{lem:5c_target}
Restrict scope to substrate-targeting adversarial events
(Definition~\ref{def:shock}).  Under invariant $I_6'$ ($\mindep \geq
m^* \geq 3$), Paper~3's anti-monopolar conclusion holds in a static
safe region strictly larger than the linearized form admits.
Specifically, diversity dominates whenever
\begin{equation}
\label{eq:safe_region}
\Delta r_K \;<\; \rext + (1 - \gamma)\bigl(\Ddiv - \Delta_0\bigr),
\end{equation}
where $\Ddiv$ is the substrate-diversification value advantage:
\begin{equation}
\label{eq:ddiv_definition}
\Ddiv \;=\; \volL_K \cdot \bigl(\lossDmax - \lossdivmax\bigr) \cdot
\mathbb{E}\!\left[\gamma^{\Tadv}\right],
\end{equation}
$\lossDmax$ is the dominator's worst-case shock-loss fraction,
$\lossdivmax$ is the diversity strategy's worst-case shock-loss
fraction, and $\mathbb{E}[\gamma^{\Tadv}]$ is the expectation of the
discount factor over the adversarial event time distribution.
\end{lemma}

The structural claim: under substrate-targeting adversarial events
with positive $\Ddiv$, diversity has a survival-value advantage that
the linearized form (\citep[Proposition~6 Corollary]{lasser2026horizon})
does not capture.  The advantage extends the safe region into part of
regime~(iii) of Paper~3's strategy-dependent corollary, conditional on
the worst-case loss-fraction inequality $\lossDmax > \lossdivmax$
holding.

\section{The model extension}

Paper~3's linearized argument
(\citep[Proposition~6]{lasser2026horizon}) compares two infinite-horizon
strategies under deterministic growth, with no stochastic adversarial
events.  Lemma~5c adds substrate-targeting shocks to the model.

\begin{definition}[Substrate-Targeting Shock]
\label{def:shock}
A \emph{substrate-targeting shock} eliminates the
$\volL$-contribution of one substrate $s$ from the coalition's
capability poset, including (a)~individual capabilities supported on
$s$ and (b)~cooperative capabilities whose participants required $s$
for production.  Formally, the post-shock coalition state has
\[
  \volL_K^{\mathrm{post}}(s\text{-shock}) \;=\;
  \volL\bigl(G_K^{-s}\bigr)
\]
where $G_K^{-s}$ is the coalition's capability poset with all
capabilities supported on $s$ (individual and cooperative)
removed.

\textbf{Joint / event-class independence assumption.}  No single
adversarial mechanism in scope produces shocks to more than one
substrate.  Formally: for each adversarial mechanism class
$\mathcal{A}$ in the deployment's threat model, the support of any
shock event drawn from $\mathcal{A}$ is contained in a single
substrate $s$.  This is strictly stronger than pairwise
failure-correlation independence (Definition~\ref{def:indep}), which
would only require the joint failure distribution to factorize across
substrate pairs.

Shock arrivals form a marked point process with marks indicating the
targeted substrate.  We write $\Tadv$ for the time of the first shock,
treated as a random variable with finite expectation.
\end{definition}

\begin{remark}[Why this scope is restrictive]
The substrate-targeting shock is the natural adversarial event for
which substrate diversity supplies the survival advantage. Other
adversarial classes do not distinguish diversity from domination at
the substrate level and produce different (or zero) $\Ddiv$:
\begin{itemize}
\item \emph{Capability-targeting shocks} suppress individual
  capabilities regardless of their substrate placement; diversity
  doesn't reduce the per-capability shock probability.
\item \emph{Coalition-internal corruption} compromises the actor's
  reasoning or the agent population's coordination; substrate
  diversity does not protect against an insider that operates within
  every substrate.
\item \emph{Environment manipulation} (Lemma~5e's domain) alters true
  $\rext$ or $\Delta r_K$ via exogenous variables; the
  substrate-survival mechanism does not apply.
\end{itemize}

Lemma~5c therefore protects against substrate-targeting attackers
specifically: physical infrastructure attack, single-substrate
compromise (e.g., shared-vulnerability LLM jailbreak across all
silicon agents on one substrate), substrate-correlated failure
modes.  It does not by itself protect against the other classes.
A deployment claim covering broader threat classes requires a
unified adversarial-event model with separate shock-class terms; this
is open work for Paper~10's main theorem composition.
\end{remark}

\section{Defining $\Ddiv$ via Paper 3's volume-survival penalty}

The v1 draft modeled $\Ddiv$ as a discounted infinite tail of
post-shock rates.  Codex Round~A found this conflated Paper~3's
one-time minimax contraction with a continuous-rate model not
derivable from Paper~2's axioms.  The v2 derivation uses Paper~3's
actual volume-survival penalty, which exists in source-paper machinery
at \citep[\S domination\_kills]{lasser2026horizon}.

\begin{definition}[Worst-Case Shock-Loss Fraction --- counterfactual form]
\label{def:loss_fraction}
For a coalition state $G_K$ with capability poset and
$\volL$-distribution across substrates $S$, the \emph{shock-loss
fraction} for substrate $s \in S$ is the counterfactual loss-fraction
under removal of $s$:
\[
  \alpha_s(G_K) \;=\; \frac{\volL(G_K) - \volL(G_K^{-s})}{\volL(G_K)},
\]
where $G_K^{-s}$ is the coalition's capability poset with all
capabilities supported on $s$ (individual and cooperative) removed.

The \emph{worst-case shock-loss fraction} under adversarial targeting
is
\[
  \ell^{\max}(G_K) \;=\; \max_{s \in S} \alpha_s(G_K).
\]
\end{definition}

\begin{remark}[Why counterfactual, not additive --- codex Round B Q1]
\label{rem:counterfactual_alpha}
An earlier draft described $\alpha_s$ as an additive distribution
across substrates, with cooperative weight ``incremented'' into each
participating substrate's share.  Codex Round~B flagged this as
unsafe: with multi-substrate cooperatives, additive attribution can
overcount.  A three-way cooperative whose participants are all three
substrates would be added to all three $\alpha_s$ terms, making them
sum above 1.

The counterfactual definition above avoids this: each $\alpha_s$ is
an independent counterfactual under removal of $s$, not a partition
share.  The shock-loss profile $(\alpha_1, \ldots, \alpha_m)$ may
sum above 1 when there are cross-substrate cooperative
dependencies, but each individual $\alpha_s \in [0, 1]$.

Additive attribution remains useful as an \emph{estimator} of $\alpha_s$
when the counterfactual computation is intractable, but it is an
estimator with known overcounting bias when downward-closure overlap
is ignored.  Paper 10's deployment-tooling section should specify the
counterfactual form as the operational target and the additive form
only as a fallback approximation with stated correction.
\end{remark}

For the diversity strategy under invariant $I_6'$, balanced
$\volL$-distribution gives $\lossdivmax \approx 1/\mindep$.  For
$\mindep = 3$, $\lossdivmax \approx 1/3$.

For the dominator strategy under full domination on a single
failure-correlated substrate, $\lossDmax \approx 1$ (the entire
coalition $\volL$ is supported on the dominant substrate).  Under
partial domination with residual operations on other substrates,
$\lossDmax$ is bounded above by $1$ but could be smaller.

\begin{definition}[Substrate-Diversification Advantage $\Ddiv$]
\label{def:ddiv}
Under the substrate-targeting shock model
(Definition~\ref{def:shock}) and worst-case adversarial substrate
selection,
\[
  \Ddiv \;=\; \volL_K \cdot \bigl(\lossDmax - \lossdivmax\bigr)
  \cdot \mathbb{E}\!\left[\gamma^{\Tadv}\right].
\]
Each factor:
\begin{itemize}
\item $\volL_K$ scales the absolute $\volL$ at risk to the
  pre-event coalition volume.
\item $\lossDmax - \lossdivmax$ is the worst-case loss-fraction
  advantage of diversity over domination.  Under full domination on
  a failure-correlated substrate with $\mindep \geq 3$ and balanced
  diversity, this is approximately $1 - 1/\mindep$.
\item $\mathbb{E}[\gamma^{\Tadv}]$ is the expectation of the discount
  factor over the random shock arrival time, \emph{not}
  $\gamma^{\mathbb{E}[\Tadv]}$ (Paper~3 \S risk.tex:60 explicitly
  warns about this conflation; heavy-tailed $\Tadv$ produces
  systematically different values for the two).
\end{itemize}
\end{definition}

\paragraph{Sign condition.}
Lemma~5c's safe region tightening
(Equation~\ref{eq:safe_region}) is positive only when
$\lossDmax > \lossdivmax$.  In the balanced full-domination case,
this is $1 > 1/\mindep$, which holds for any $\mindep > 1$.  In
partial-domination or non-balanced cases, the sign condition must be
checked: an unbalanced diversity strategy with one substrate carrying
most volume can have $\lossdivmax \approx \lossDmax$, eliminating the
advantage.

\section{Substrate-distinctness via $\mindep$ and the tripartite
canonical identification}

Lemma~5c's $\mindep$ counts substrates at the failure-correlation
independence level, not nominal substrate type.  This excludes
``skeleton substrates'' that nominally separate but share failure
modes (e.g., three silicon agents from different vendors but
identical training data; two cloud deployments on different providers
but same operating substrate; etc.).

\begin{definition}[Joint / Event-Class Failure-Correlation Independence]
\label{def:indep}
Substrates $s_1, \ldots, s_k$ are \emph{jointly failure-correlation
independent in the event-class sense} relative to a deployment's
threat model $\mathcal{T}$ if: for every adversarial mechanism class
$\mathcal{A} \in \mathcal{T}$, the support of any shock event drawn
from $\mathcal{A}$ is contained in a single substrate.

This is strictly stronger than pairwise correlation factorization
(which only requires pair-events to factor) and matches the
``no shock spans multiple substrates'' condition the safe-region
calculation actually needs.
\end{definition}

\begin{remark}[Mechanism = causal campaign, not atomic step --- codex Round B Q4]
\label{rem:mechanism_campaign}
``Adversarial mechanism'' in Definition~\ref{def:indep} must refer to
the whole causal attack episode or campaign, not a single atomic step
within it.  An attacker can construct multi-step campaigns where each
step targets one substrate (satisfying step-level independence) but
the aggregate campaign hits multiple substrates by design.  Treating
each step as a separate mechanism would falsely satisfy
Definition~\ref{def:indep} while the aggregate violation is real.

Operational implication: the threat model $\mathcal{T}$ must enumerate
attack \emph{campaigns} (or causal episodes that produce shock events
across the deployment lifetime) and verify each campaign's support
is contained in a single substrate.  A causal-graph or
campaign-level factorization would compose more cleanly with Paper
10's other invariants; v3 keeps the event-class formulation but
explicitly requires campaign-level interpretation.

Unknown attack mechanisms not in $\mathcal{T}$ do not retroactively
break the math --- they break the assumption that the threat model is
complete.  This is the standard threat-model-completeness residual
inherited from any security analysis.
\end{remark}

\begin{definition}[$\mindep$]
\label{def:mindep}
$\mindep$ is the maximum cardinality of a subset of nominal substrates
that is jointly failure-correlation independent in the event-class
sense (Definition~\ref{def:indep}).
\end{definition}

\subsection{The canonical tripartite identification}
\label{sec:tripartite}

Per Paper~10 \S6, the canonical three-substrate identification meeting
$\mindep \geq 3$ is:

\begin{itemize}
\item \textbf{Substrate 1 (Human):} judgment, strategic decisions,
  value reasoning, governance vote.
\item \textbf{Substrate 2 (AI):} planning, generation, execution,
  broad-recall, computation.
\item \textbf{Substrate 3 (Formal-Operational):} verification,
  attestation, commitment, audit, monitoring, governance enforcement.
\end{itemize}

These are joint / event-class failure-correlation independent: prompt
injection breaks neither human judgment nor cryptographic commitments;
cognitive bias breaks neither LLM inference nor ledger integrity;
trusted-setup failure breaks neither human reasoning nor LLM behavior.

In contrast, a ``Claude + Codex + human'' identification fails the
joint independence test because the two LLMs share substantial
adversarial-event surface (transferable prompt injections, training
data overlap, architectural similarity).  Operationally,
$\mindep = 2$ for that configuration.

\subsection{Skeleton substrates and the audit burden}

If a deployment uses $\meff \geq 3$ but $\mindep < 3$, an adversary
with a single attack mechanism can take down multiple ``substrates''
simultaneously, and the safe region calculation does not apply.  This
is the codex-flagged failure mode.

Operationally, $\mindep$ is determined by failure-mode auditing
across the substrate population.  Required: enumerate the threat
model $\mathcal{T}$, list the adversarial mechanism classes, verify
that each is contained within a single substrate.  This is
non-trivial in practice (the audit may identify previously-unknown
shared vulnerabilities) but tractable for well-understood substrate
configurations.

\section{Adversarial targeting and the balanced-$\volL$ requirement}

The v1 draft assumed shocks select substrates independently of
$\volL$-distribution.  Under adversarial targeting, the attacker
chooses the substrate to maximize $\volL$-loss, so the worst-case
analysis must use $\max_s \alpha_s$ rather than a random average.

\subsection{The corrected balance threshold}

The codex-corrected balanced-$\volL$ requirement is:
\begin{equation}
\label{eq:balance}
\max_s \alpha_s \;\leq\; \frac{1}{m^*} + \epsilon
\end{equation}
for small $\epsilon > 0$, where $\alpha_s$ is the full shock-loss
fraction (Definition~\ref{def:loss_fraction}) and $m^* = \mindep$.

\textbf{The v1 error.}  v1 stated this threshold as $\max_s \alpha_s
\leq 1 - 1/m^* + \epsilon$.  For $m^* = 3$, that allows
$\alpha_{\max} = 2/3 + \epsilon$ --- two-thirds of value on a single
substrate, which is the opposite of balance.  The corrected form
$\alpha_{\max} \leq 1/3 + \epsilon$ requires substrates to be
near-uniform, which is what ``balance'' actually means.

The intuition: ``balanced'' means \emph{no} substrate has an
overwhelming share, which is the \emph{low max-share} regime, not the
high one.  Sanity check at $m^* = 3$: balanced requires
$\alpha_{\max} \approx 1/3$; the v1 threshold permitted
$\alpha_{\max}$ up to $2/3$, which is monopoly-ish.

\subsection{Operational implication: $\alpha$-distribution audit}

The deployment-tooling section of Paper~10 must include an
$\alpha$-distribution audit: for each substrate, estimate its full
shock-loss fraction (individual + cooperative-loss attribution) and
verify $\max_s \alpha_s \leq 1/m^* + \epsilon$.

The cooperative-loss attribution is the subtle part.  If the system's
high-leverage cooperative requires participation from all three
substrates, removing any one substrate eliminates the cooperative
entirely; each substrate's $\alpha$ then includes a share of the
cooperative's value.  In the worst case, three-way cooperatives
inflate each substrate's $\alpha$ to roughly $\frac{1}{3}$ (individual)
$+\frac{1}{3}$ (cooperative) $= \frac{2}{3}$, blowing past the
balance threshold.

\textbf{Design implication:} cooperative capabilities should
\emph{not} all depend on every substrate.  Either distribute
cooperatives so multiple cooperatives can substitute for a lost one,
or accept that the system has a higher effective $\alpha_{\max}$ and
the safe region is correspondingly narrower.  This is a real
deployment-design constraint Paper~10 must call out.

\section{Composing with the linearized form}

Lemma~\ref{lem:5c_target}'s safe region
(Equation~\ref{eq:safe_region}) extends Paper~3's linearized region
($\Delta r_K < \rext$) by the term $(1-\gamma)(\Ddiv - \Delta_0)$.

\begin{itemize}
\item \textbf{Static $\Delta_0 \leq 0$ (Paper~3 Case~1).}  Immediate
  $\volL$-change from domination is non-positive; both $-\Delta_0$
  and $\Ddiv$ contribute positively to the safe region.  The added
  $\Ddiv$ term is pure tightening.
\item \textbf{Static $\Delta_0 > 0$ (Paper~3 Case~2).}  The $\Ddiv$
  term competes with the $-\Delta_0$ term.  Net contribution is
  positive when $\Ddiv > \Delta_0$, i.e., when the
  substrate-survival advantage exceeds the immediate domination
  bonus.  This is the practically interesting case.
\end{itemize}

Substituting the explicit form:
\begin{equation}
\label{eq:safe_region_concrete}
\Delta r_K \;<\; \rext + (1 - \gamma)\bigl[\volL_K (\lossDmax -
\lossdivmax) \mathbb{E}[\gamma^{\Tadv}] - \Delta_0\bigr].
\end{equation}

This is the operationally checkable safe region.  Each quantity is
either invariant-bounded or empirically estimable:
\begin{itemize}
\item $\rext$: bounded below via Lemma~5a substrate-floor argument.
\item $\Delta r_K$: compared against the bound; requires upper-bound
  estimate from the deployment's adversarial threat model.
\item $\volL_K$: ledger-observable from the coalition's accumulated
  capability volume.
\item $\lossDmax$: under full domination on a single failure-
  correlated substrate, $\approx 1$; partial domination requires
  specific estimation.
\item $\lossdivmax$: bounded by $1/m^* + \epsilon$ under
  $\alpha$-distribution audit; otherwise empirically the maximum
  shock-loss fraction across substrates.
\item $\mathbb{E}[\gamma^{\Tadv}]$: empirically estimated from
  historical shock-rate data and the deployment's $\gamma$.
\item $\Delta_0$: computable from Paper~3's immediate
  $\volL$-change formula given the coalition state.
\end{itemize}

\section{Multi-shock extension via discounted marginal loss increments}

The single-shock formula
(Equation~\ref{eq:safe_region_concrete}) is valid when shock arrival
is rare relative to the planning horizon ($\mathbb{E}[\Tadv]$
comparable to or longer than the discounted planning window).  For
high-$\pshock$ regimes, multiple shocks accumulate and the bound must
account for the cumulative effect.

\subsection{Why the survival-state formulation fails}
\label{sec:why_survival_fails}

An earlier attempt (v2) defined a survival-state stream
$A_{\mathrm{multi}} = \mathbb{E}\sum_t \gamma^t (\volL_{\mathrm{div}}^{\mathrm{surv}}(t) - \volL_D^{\mathrm{surv}}(t))$.
Codex Round~B flagged a dimensional inconsistency: in the single-shock
limit, the surviving-volume difference \emph{persists indefinitely}
after the shock, so the discounted stream evaluates to $\Ddiv / (1 -
\gamma)$ rather than $\Ddiv$.  Substituting the stream into
Equation~\ref{eq:safe_region} therefore changes the dimensional scale.

Additionally, v2's $(1 - \alpha_{\max})^k$ multiplicative survival
formula was wrong: with $m$ distinct equal substrates removed without
replacement, surviving fraction is additive ($1 - k/m$), not
multiplicative.  The multiplicative form would describe repeated
proportional hits to the remaining state, not substrate exhaustion.

The v3 formulation corrects both errors by using discounted marginal
loss increments at shock times.

\subsection{The discounted-marginal-loss formulation}
\label{sec:multishock_marginal}

\begin{definition}[Multi-Shock Diversification Advantage]
\label{def:ddiv_multi}
Let $\{T_i\}_{i=1}^{\infty}$ denote the (random) shock arrival times,
with $T_i < T_{i+1}$.  Let $\delta_i^{\mathrm{div}}, \delta_i^D$
denote the marginal $\volL$-loss for the diversity and domination
strategies at shock $i$, respectively.  The
\emph{multi-shock substrate-diversification advantage} is the
discounted sum of marginal loss differences:
\begin{equation}
\label{eq:ddiv_multi}
\Ddiv^{\mathrm{multi}}
  \;=\; \mathbb{E}\!\left[\sum_{i=1}^{\infty}
  \gamma^{T_i} \cdot
  \bigl(\delta_i^D - \delta_i^{\mathrm{div}}\bigr)\right].
\end{equation}
Note: $\delta_i^D - \delta_i^{\mathrm{div}}$ is positive when the
dominator loses more than diversity at shock $i$, negative when
diversity loses more (the dominator is already at zero on subsequent
shocks).  The expectation is over the joint distribution of
$\{T_i\}$ and the substrate-targeting policy.
\end{definition}

\subsection{Single-shock reduction (sanity check)}
\label{sec:multishock_marginal}

In the single-shock model (where $N = 1$ a.s., i.e., the shock
process produces exactly one shock event with probability 1, and
$T_i = \infty$ for $i \geq 2$ with $\gamma^{\infty} = 0$ by
convention), $\Ddiv^{\mathrm{multi}}$ reduces to:
\[
  \Ddiv^{\mathrm{multi}} \;\to\; \mathbb{E}[\gamma^{\Tadv}] \cdot
  (\delta_1^D - \delta_1^{\mathrm{div}}).
\]
At the first shock, under full failure-correlated domination on a
single substrate, the dominator loses everything: $\delta_1^D = \volL_K$.
Under balanced diversity with $\mindep \geq 3$, the diversity
strategy loses $\lossdivmax \cdot \volL_K \approx \volL_K / \mindep$.
The advantage is:
\[
  \delta_1^D - \delta_1^{\mathrm{div}} = \volL_K - \lossdivmax \volL_K
  = \volL_K(1 - \lossdivmax) = \volL_K(\lossDmax - \lossdivmax)
\]
(using $\lossDmax = 1$ under full domination).

This recovers the single-shock $\Ddiv$ from
Definition~\ref{def:ddiv} exactly.  The dimensional consistency
that v2's $A_{\mathrm{multi}}$ lost is restored.

\paragraph{Note on single-shock vs $\pshock \to 0$ limit.}
The single-shock model ($N=1$ a.s.) is a different formal object
from a Poisson model with $\lambda \to 0$ on infinite horizon.
Under Poisson arrivals on an infinite horizon, later shocks still
occur a.s.; their discounted contribution vanishes as
$\kappa = \lambda/(\lambda - \ln \gamma) \to 0$ but the formal
expectation is over the Poisson process, not a single random time.
The two models agree numerically in the $\kappa \to 0$ limit but
are not the same statistical structure.  Paper~10's deployment
analysis should be explicit about which model is in force; the
single-shock $\Ddiv$ form (Definition~\ref{def:ddiv}) uses the
$N=1$ a.s. model directly with $\Tadv$ as the random shock time.

\subsection{Substrate exhaustion under additive removal --- balanced case}
\label{sec:substrate_exhaustion}

Codex Round~B noted that $(1 - \alpha_{\max})^k$ is wrong for
substrate exhaustion: with $m$ distinct equal substrates removed
without replacement, surviving fraction is additive ($1 - k/m$), not
multiplicative.  The corrected accounting holds for the
\textbf{balanced, additive-cooperative special case}; the general case
requires sequential counterfactuals (\S\ref{sec:sequential_counterfactual}).

Under independent substrate-targeting shocks where each shock removes
one previously-functional substrate (no replacement), and assuming
the substrate-cooperative structure is dominated by cross-substrate
cooperatives rather than redundant capabilities (the
cooperative-overlap regime), the diversity strategy's marginal loss
at shock $i$ (for $i \leq \mindep$, balanced case) is:
\[
  \delta_i^{\mathrm{div}} \;\leq\; \volL_K / \mindep
  \quad\text{(approximate equality in balanced case;
  upper bound under cooperative-overlap)}.
\]
The dominator's marginal loss after shock 1 is zero (already at
$\volL = 0$): $\delta_i^D = 0$ for $i \geq 2$.

Substituting into Equation~\ref{eq:ddiv_multi}:
\begin{align*}
\Ddiv^{\mathrm{multi}}
  &\;\geq\; \mathbb{E}[\gamma^{T_1}] \cdot \volL_K (1 - 1/\mindep)
   - \sum_{i=2}^{\mindep} \mathbb{E}[\gamma^{T_i}] \cdot \volL_K /
   \mindep \\
  &\;=\; \volL_K \left[\mathbb{E}[\gamma^{T_1}] (1 - 1/\mindep)
   - \frac{1}{\mindep} \sum_{i=2}^{\mindep} \mathbb{E}[\gamma^{T_i}]\right].
\end{align*}
The second sum is finite (bounded by $(\mindep - 1)$ terms each
$\leq 1$) and represents the erosion of diversity's advantage as it
also exhausts substrates.

The inequality direction (lower bound on $\Ddiv^{\mathrm{multi}}$)
follows from the cooperative-overlap regime assumption: equal
analysis overcounts diversity's later marginal losses (because some
cooperatives are already lost from earlier shocks), so the
formula above is a conservative lower bound on the actual
sequential-counterfactual $\Ddiv^{\mathrm{multi}}$.  This is what
deployment-safety claims need.

\subsection{The general case: sequential counterfactuals}
\label{sec:sequential_counterfactual}

For deployments outside the cooperative-overlap regime --- in
particular, deployments with substantial redundant-capability
structure (capabilities replicated across multiple substrates) ---
the equal-substrate analysis using original-counterfactual $\alpha_s$
is anti-conservative.  Codex Round~C identified this:

\begin{quote}
``v3's $\alpha_s$ is a single-removal counterfactual, not an additive
partition. Sequential losses must be computed as
$\delta_i^{\mathrm{div}} = \volL(G^{-R_{i-1}}) - \volL(G^{-(R_{i-1}
\cup \{s_i\})})$ where $R_{i-1}$ is the set already removed.
Overlapping cooperatives can make original $\alpha_s$ overcount later
losses; redundant capabilities can make original $\alpha_s$
undercount later losses after redundancy is exhausted.''
\end{quote}

In the redundancy regime, equal analysis underestimates $\delta_i$
for late shocks (the cumulative loss appears smaller than actual),
which overestimates $\Ddiv^{\mathrm{multi}}$, which gives an
anti-conservative (too-permissive) safe-region inequality.  Lemma~5c
v4 explicitly does \emph{not} cover this regime.

\paragraph{Cooperative-vs-redundancy audit.}
For a deployment to invoke Lemma~5c's multi-shock result, the
substrate-cooperative structure must be dominated by cross-substrate
cooperatives.  Operationally, this requires verifying:
\begin{enumerate}
\item For each capability $c$ in the poset, identify the substrate(s)
  it is supported on.
\item Categorize: (a)~single-substrate capabilities,
  (b)~cross-substrate cooperatives (one capability whose production
  requires participation from multiple substrates),
  (c)~redundant capabilities (capability replicated identically on
  multiple substrates).
\item Verify category (b) volume exceeds category (c) volume in
  $\volL$-weighted sum.
\end{enumerate}
The canonical tripartite identification (Human / AI /
Formal-Operational) passes this audit by construction: the
high-leverage capabilities (verified workflow, governance
enforcement) are cross-substrate cooperatives, not redundant
capabilities.

\paragraph{Sequential counterfactual derivation as open work.}
For deployments where redundancy dominates or the regime is mixed,
the multi-shock result requires deriving $\delta_i$ from sequential
counterfactuals directly:
\[
  \delta_i^{\mathrm{div}}(R_{i-1}) \;=\; \volL\bigl(G_K^{-R_{i-1}}\bigr)
  - \volL\bigl(G_K^{-(R_{i-1} \cup \{s_i\})}\bigr).
\]
This requires evaluating $\volL$ on $2^{\mindep}$ counterfactual
poset states, which is computationally tractable for $\mindep = 3$
but grows combinatorially.  Paper~10 cites the sequential
counterfactual extension as Lemma~5c-prime (open work) for
deployments outside the cooperative-overlap regime.

\subsection{High-$\lambda$ behavior (continuous-time Poisson model)}

We adopt a \textbf{continuous-time homogeneous Poisson model} for
shock arrivals: shocks arrive at rate $\lambda$ per unit time
($T_i - T_{i-1} \sim \mathrm{Exp}(\lambda)$ iid).  Codex Round~C
flagged that earlier drafts using ``$\pshock$ per period'' was
ambiguous: a discrete-time per-period probability model would change
the $\kappa$ formula.  V4 standardizes on the continuous-time
$\lambda$ throughout this section.

For Poisson shock arrivals with rate $\lambda$,
$\mathbb{E}[\gamma^{T_i}] = \kappa^i$ where
\[
  \kappa \;=\; \mathbb{E}[\gamma^{T_1}] \;=\;
  \frac{\lambda}{\lambda - \ln \gamma}.
\]

\paragraph{Discrete-time alternative.}
If a deployment uses a per-period discrete-time shock model with
per-period probability $\pshock$, then $T_1 \sim \mathrm{Geometric}(\pshock)$
and $\mathbb{E}[\gamma^{T_1}] = \pshock \gamma / (1 - (1-\pshock)\gamma)$,
which gives a different $\kappa$ formula.  Paper~10 must specify
which model the deployment uses; the continuous-time form is the
standard for adversarial-event analysis and is what v4 derives below.

Substituting into Equation~\ref{eq:ddiv_multi} (cooperative-overlap
regime, balanced case):
\[
  \Ddiv^{\mathrm{multi}}
  \;\geq\; \volL_K \left[\kappa (1 - 1/\mindep)
   - \frac{1}{\mindep} \sum_{i=2}^{\mindep} \kappa^i\right]
  \;=\; \volL_K \left[\kappa - \frac{1}{\mindep}
   \cdot \frac{\kappa(1 - \kappa^{\mindep})}{1 - \kappa}\right]
  \quad (\kappa < 1).
\]

\paragraph{Removable singularity at $\kappa = 1$.}
The closed form has a $1/(1 - \kappa)$ factor, but the singularity at
$\kappa = 1$ is \emph{removable}: the numerator
$\kappa(1 - \kappa^{\mindep})$ also vanishes at $\kappa = 1$, by
L'H\^opital's rule:
\[
  \lim_{\kappa \to 1^{-}} \frac{\kappa(1 - \kappa^{\mindep})}{1 - \kappa}
  \;=\; \mindep.
\]
Therefore $\Ddiv^{\mathrm{multi}} / \volL_K \to 1 - 1 = 0^+$ as
$\kappa \to 1$.  The advantage approaches zero from above, not
infinity.

\textbf{Operational regime condition.}  The safe region is
non-trivial only when $\Ddiv^{\mathrm{multi}}$ is bounded away from
zero.  This requires $\kappa$ bounded away from $1$:
$\kappa \leq 1 - \delta$ for some operationally-determined
$\delta > 0$.  Equivalently:
\[
  \frac{\lambda}{\lambda - \ln \gamma} \;\leq\; 1 - \delta
  \quad\Leftrightarrow\quad
  \lambda \;\leq\; \frac{(1-\delta)(- \ln \gamma)}{\delta}.
\]

Paper~10's deployment-tooling section should state the regime
condition on $\kappa$, not on $\lambda$ or $\pshock$ directly:
shocks arrive slowly enough relative to discounting that
$\kappa < 1 - \delta$.  For typical $\gamma$ near $1$ ($\ln \gamma$
small and negative), this requires $\lambda$ small (rare shocks).
For $\gamma$ farther from $1$ (high time preference, short effective
horizon), the regime condition is less restrictive on $\lambda$.

\textbf{What this means operationally.}  When shocks are very
frequent relative to discounting ($\kappa$ close to $1$), the
safe-region advantage from substrate diversity becomes trivially
small even with all invariants satisfied.  This is not a failure of
the lemma; it is a correct statement of the regime where substrate
diversification is and isn't valuable.  Deployments operating in
high-$\lambda$ environments (e.g., adversarial environments with
frequent attacks) need additional defenses beyond substrate
diversity to invoke the deployment-safety theorem.

\subsection{Replacement in the safe-region inequality}

For the safe-region inequality, replace $\Ddiv$ in
Equation~\ref{eq:safe_region} with $\Ddiv^{\mathrm{multi}}$:
\[
  \Delta r_K \;<\; \rext + (1 - \gamma)\bigl(\Ddiv^{\mathrm{multi}}
  - \Delta_0\bigr).
\]
The dimensional consistency is preserved because
$\Ddiv^{\mathrm{multi}}$ is a sum of $\Ddiv$-style penalties (each a
one-time loss difference scaled by an expected discount factor), not
an integrated survival stream.

\section{Substrate-targeting scope (codex Q5)}

Lemma~5c is explicitly scoped to substrate-targeting adversaries.
Other adversarial classes are not covered:

\begin{itemize}
\item \textbf{Capability-targeting:} an adversary suppresses an
  individual capability $c$ regardless of substrate.  Substrate
  diversity does not reduce the per-capability shock probability;
  the safe-region argument does not apply.
\item \textbf{Coalition-internal corruption:} an adversary
  compromises the actor's reasoning, world model, or coordination
  protocol.  This is Paper~6's domain (Lyapunov degradation,
  cascade dynamics) rather than substrate diversification.
\item \textbf{Environment manipulation:} an adversary alters true
  $\rext$ or $\Delta r_K$ via exogenous environment variables.
  Lemma~5e and invariant $I_8$ (environment-side witnesses) handle
  this; substrate diversification does not.
\item \textbf{Coordinated multi-substrate attack:} a single
  adversarial mechanism that produces shocks to multiple
  substrates simultaneously.  Excluded by joint / event-class
  independence (Definition~\ref{def:indep}); if the threat model
  contains such mechanisms, $\mindep$ collapses and the safe
  region does not apply.
\end{itemize}

For Paper~10's deployment-safety theorem to cover broad adversarial
threat models, additional shock-class-specific terms are needed.
These are open work; Lemma~5c covers only the substrate-targeting
class.

\paragraph{Citation discipline (codex Round B Q5).}
Lemma~5c should be cited only for substrate-targeting survival
advantage.  Paper~10 \S6's broader claims about cooperative anchoring
do cover capture-of-existing, asymmetric capture, cooperative
forking, and time-asymmetry capture, but those rely on separate
operational invariants ($I_9$ substrate-exclusivity observability,
$I_{10}$ coverage/materiality, $I_{11}$ latency bounds) plus Paper~5's
structural defenses against governance/coalition capture.  Lemma~5c
does \emph{not} establish those broader claims.  When Paper~10 cites
Lemma~5c in the deployment-safety theorem composition, the citation
should be scoped: ``for substrate-targeting adversarial events, the
diversity-domination value comparison admits the safe region of
Equation~\ref{eq:safe_region_concrete},'' not ``Lemma~5c establishes
the deployment claim.''

\section{Honest residuals}

\paragraph{Paper~3's volume-survival form requires explicit identification.}
Lemma~5c uses $\Ddiv = \volL_K \cdot (\lossDmax - \lossdivmax) \cdot
\mathbb{E}[\gamma^{\Tadv}]$, derived from Paper~3's
\citep[\S domination\_kills]{lasser2026horizon} machinery.  Paper~3
itself states the contraction in absolute-volume form
($\Ddiv = \min_s \volL(G^{-s})$, codex Round~A pointer
substitution.tex:1381); the normalized loss-fraction form here is a
direct rewriting under the $\alpha_s$ decomposition, but Paper~10
should cite the equivalence carefully.

\paragraph{$\pshock$ is empirically estimated.}
The per-period shock probability is operational input.  An operator
that under-estimates $\pshock$ over-estimates the safe region.
Paper~10's deployment-tooling section must specify how $\pshock$ is
computed (worst-case from shock taxonomy?  Bayesian update?
Governance-set floor?).

\paragraph{Cooperative-loss attribution well-defined under counterfactual form.}
Under v3's counterfactual definition of $\alpha_s$
(Definition~\ref{def:loss_fraction}), cooperative-loss attribution
is unambiguous: $\alpha_s$ is computed from the single counterfactual
$\volL(G_K^{-s})$ (substrate $s$ removed), without splitting
cooperative weight across substrates.  The shock-loss profile
$(\alpha_1, \ldots, \alpha_m)$ may sum above 1 with cross-substrate
cooperative dependencies, which is correct: removing different
substrates loses different (overlapping) sets of cooperatives, and
each counterfactual stands alone.

The earlier additive attribution (which would have required canonical
rules) is now an estimator, not the definition
(Remark~\ref{rem:counterfactual_alpha}).

\paragraph{Joint independence is hard to verify.}
Definition~\ref{def:indep}'s ``no single adversarial mechanism in
scope spans multiple substrates'' is a strong condition.  In practice,
it requires enumerating the threat model's adversarial mechanism
classes and verifying each is contained within one substrate.  Novel
attack mechanisms not in the original threat model can violate the
condition retroactively, invalidating the safe region.

\paragraph{Partial-domination case requires care.}
Under partial domination with residual operations on other
substrates, $\lossDmax < 1$.  The sign condition $\lossDmax >
\lossdivmax$ may fail in some partial-domination configurations,
collapsing the safe region.  Paper~10 should characterize the
partial-domination regimes where Lemma~5c applies.

\paragraph{Multi-shock formulation assumes no recovery.}
Definition~\ref{def:ddiv_multi} treats shocks as independent and
substrate-removal as permanent.  Real deployments may include
substrate-recovery dynamics (failed substrates can be repaired or
replaced); multi-shock with recovery is its own analysis with a
different functional form for $\delta_i$ and is not derived here.
Under recovery, the diversity strategy's substrate-exhaustion
trajectory is bounded above by $\mindep$ but may stabilize at a
non-zero count, changing the high-$\pshock$ asymptotic.

\paragraph{Tripartite identification scope.}
The canonical identification (Human / AI / Formal-Operational)
satisfies $\mindep = 3$ under standard threat models, but novel
threat models could correlate them (e.g., a cyber-physical attack
that compromises both formal infrastructure and human operators
simultaneously, or social-engineering that affects human judgment via
AI-generated misinformation).  Paper~10 must specify which threat
models the canonical identification is valid for.

\section{Summary of changes from v3 (codex Round C)}

\begin{enumerate}
\item \textbf{Q1 (single-shock formalization):} Tightened single-shock
  reduction from ``$T_i = \infty$ for $i \geq 2$ with negligible
  probability'' to formal $N = 1$ a.s. with $\gamma^{\infty} = 0$
  convention (\S\ref{sec:multishock_marginal}).  Added a paragraph
  distinguishing the $N=1$ a.s. model from the $\lambda \to 0$
  Poisson limit (different formal objects, agree numerically only
  in the $\kappa \to 0$ limit).

\item \textbf{Q2 (general non-equal scope restriction):}
  Multi-shock claim explicitly scope-restricted to the
  cooperative-overlap regime where equal-substrate analysis is a
  conservative lower bound on $\Ddiv^{\mathrm{multi}}$.  The
  redundancy regime requires sequential counterfactual derivation
  (\S\ref{sec:sequential_counterfactual}) and is named as
  ``Lemma~5c-prime'' for follow-up.  Added cooperative-vs-redundancy
  audit as an operational requirement.

\item \textbf{Q3 (Poisson rate vs per-period probability):} $\pshock$
  reframed as continuous-time Poisson rate $\lambda$.  Discrete-time
  per-period probability model produces a different $\kappa$ formula
  (noted explicitly).  V4 standardizes on continuous-time as the
  default.

\item \textbf{Q3 (removable singularity at $\kappa = 1$):}  The
  closed form has a removable singularity at $\kappa = 1$:
  $\Ddiv^{\mathrm{multi}}/\volL_K \to 0^+$ via L'H\^opital, not
  $\infty$.  Operational regime condition is on $\kappa$ bounded
  away from 1 (not $\lambda$ directly).

\item \textbf{Abstract regression fix:}  v3 abstract referenced the
  v2 survival-stream $A_{\mathrm{multi}}$ formulation.  V4 abstract
  fully reflects the v3 marginal-loss formulation and the v4 scope
  restriction.
\end{enumerate}

\section{Summary of changes from v2 (codex Round B)}

\begin{enumerate}
\item \textbf{Q1 ($\alpha_s$ definitional caveat):} Replaced
  additive-distribution language with the counterfactual formula
  $\alpha_s = (\volL(G_K) - \volL(G_K^{-s})) / \volL(G_K)$
  (Definition~\ref{def:loss_fraction}).  Additive attribution
  remains useful as an estimator with stated overcounting bias
  (Remark~\ref{rem:counterfactual_alpha}).

\item \textbf{Q3 (multi-shock dimensional inconsistency):} Replaced
  the survival-state stream $A_{\mathrm{multi}}$ with discounted
  marginal loss increments at shock times
  (Definition~\ref{def:ddiv_multi}).  This reduces correctly to
  $\Ddiv$ in the single-shock limit
  (\S\ref{sec:multishock_marginal}).  Substrate exhaustion is
  additive ($1 - k/m$), not multiplicative
  (\S\ref{sec:substrate_exhaustion}).  High-$\pshock$ analysis
  shows the advantage does not collapse but the mechanism is
  front-loaded at the first shock, with subsequent erosion bounded.

\item \textbf{Q4 (mechanism scope):} Added clarification that
  ``adversarial mechanism'' in Definition~\ref{def:indep} refers to
  the whole causal attack campaign, not atomic steps within it
  (Remark~\ref{rem:mechanism_campaign}).

\item \textbf{Q5 (citation discipline):} Added explicit note that
  Lemma~5c covers substrate-targeting adversaries only; \S6's
  broader claims rely on separate invariants $I_9$--$I_{11}$ plus
  Paper~5 defenses (in the substrate-targeting scope section).
\end{enumerate}

\section{Summary of changes from v1 (codex Round A)}

\begin{enumerate}
\item \textbf{Q1 (post-shock rate model):} Replaced proportional
  rate model $(r_K + \rext)(1 - 1/\mindep)$ (not derivable from
  Paper~2) with Paper~3's volume-survival penalty
  ($\Ddiv = \volL_K(\lossDmax - \lossdivmax)\mathbb{E}[\gamma^{\Tadv}]$).
\item \textbf{Q2 (dominator post-shock rate):} Changed from binary
  $0$-vs-rate model to continuous loss-fraction parameter
  $\lossDmax$ (defaults to $1$ under full failure-correlated
  domination, smaller under partial domination).
\item \textbf{Q3 (balance threshold):} Corrected from
  $\max_s \alpha_s \leq 1 - 1/m^* + \epsilon$ (allowed monopoly) to
  $\max_s \alpha_s \leq 1/m^* + \epsilon$ (requires near-uniform).
  $\alpha_s$ is full shock-loss fraction including cooperative loss.
\item \textbf{Q4 (multi-shock):} Added $A_{\mathrm{multi}}$
  formulation as a remark; preserves algebraic structure with
  adjusted constants; does not collapse to zero as $\pshock \to 1$.
\item \textbf{Q5 (scope):} Explicitly narrowed to
  substrate-targeting adversaries.  Other shock classes named and
  excluded.
\item \textbf{$\mathbb{E}[\gamma^{\Tadv}]$ vs $\gamma^{\mathbb{E}[\Tadv]}$:}
  Used the expectation-of-discount form throughout
  (Paper~3 risk.tex:60).
\item \textbf{Joint vs pairwise independence:} Strengthened to joint
  / event-class independence
  (Definition~\ref{def:indep}).
\item \textbf{Single $\Ddiv$ form:} v1 alternated between
  $\Ddiv \geq \ldots$ and $\Ddiv \gamma^{\Tadv} \geq \ldots$;
  v2 fixes a single form
  (Definition~\ref{def:ddiv}).
\item \textbf{Tripartite canonical identification:} Added
  \S\ref{sec:tripartite} aligning Lemma~5c with Paper~10 \S6's
  Human / AI / Formal-Operational substrate identification.
\end{enumerate}

\section{Status: v4 ready for citation in cooperative-overlap regime}

V4 addresses every concern codex Round~C flagged on v3.  The
Round~B-verified single-shock form is unchanged (still citeable).
The Round~C-flagged multi-shock issues are addressed by:
(a)~tightening the single-shock reduction to formal $N=1$ a.s.,
(b)~scope-restricting the multi-shock claim to the
cooperative-overlap regime (where equal-substrate analysis is
conservative for advantage),
(c)~reframing $\pshock$ as Poisson rate $\lambda$,
(d)~noting the removable singularity at $\kappa = 1$,
(e)~fixing the abstract regression.

\subsection{What is citeable in v4}

\begin{itemize}
\item \textbf{Single-shock $\Ddiv$ (Definition~\ref{def:ddiv}):}
  Citeable per Round~B verdict.  Counterfactual $\alpha_s$ formula
  (Definition~\ref{def:loss_fraction}) used throughout.

\item \textbf{Balance threshold (Equation~\ref{eq:balance}):}
  Citeable per Round~B verdict.  $\epsilon$ is engineering
  tolerance bounded by the safe-region positivity requirement.

\item \textbf{Multi-shock $\Ddiv^{\mathrm{multi}}$
  (Definition~\ref{def:ddiv_multi}) in cooperative-overlap regime:}
  Citeable in v4 with explicit scope restriction.  The
  equal-substrate balanced case is a conservative lower bound on
  the actual sequential-counterfactual $\Ddiv^{\mathrm{multi}}$
  when the substrate-cooperative structure is dominated by
  cross-substrate cooperatives (the canonical tripartite case).

\item \textbf{Joint/event-class independence (Definition~\ref{def:indep}):}
  Citeable per Round~B verdict, with the campaign-vs-atomic-step
  clarification (Remark~\ref{rem:mechanism_campaign}).

\item \textbf{Tripartite identification (\S\ref{sec:tripartite}):}
  Citeable per Round~B verdict, with citation discipline note in
  the substrate-targeting scope section.

\item \textbf{High-$\lambda$ regime condition:}  Citeable in v4.
  The safe region requires $\kappa < 1 - \delta$ (operationally
  determined $\delta$); deployments outside this regime need
  additional defenses beyond substrate diversity.
\end{itemize}

\subsection{What is not citeable in v4}

\begin{itemize}
\item \textbf{Multi-shock $\Ddiv^{\mathrm{multi}}$ in
  redundancy regime:} Equal-substrate analysis is anti-conservative
  when redundant capabilities dominate cross-substrate cooperatives.
  Sequential counterfactual derivation
  (\S\ref{sec:sequential_counterfactual}) is open work,
  named as ``Lemma~5c-prime'' for follow-up.

\item \textbf{Multi-shock with substrate recovery:} Definition
  assumes no recovery.  Recovery dynamics change the
  substrate-exhaustion trajectory and require a different
  $\delta_i$ formulation.  Open work.
\end{itemize}

\subsection{Operational requirements added by v4}

For Paper~10 to cite Lemma~5c's multi-shock result, the deployment
must include:

\begin{enumerate}
\item \textbf{Cooperative-vs-redundancy audit
  (\S\ref{sec:sequential_counterfactual}):} Verify the deployment's
  $\volL$ structure is dominated by cross-substrate cooperatives, not
  redundant capabilities.  The canonical tripartite identification
  (Human / AI / Formal-Operational) passes by construction.

\item \textbf{$\kappa$-regime audit:}  Verify $\kappa = \mathbb{E}[\gamma^{T_1}]$
  is bounded away from 1 by some stated $\delta$.  This depends on
  both shock arrival rate $\lambda$ and discount factor $\gamma$;
  Paper~10 should specify how $\delta$ is set.

\item \textbf{Shock model specification:}  State explicitly whether
  the deployment uses the continuous-time Poisson rate $\lambda$
  model (default) or a discrete-time per-period probability
  $\pshock$ model (different $\kappa$ formula).
\end{enumerate}

These add to the deployment-tooling burden but are operationally
implementable for the canonical tripartite case.

\subsection{V4 confirmation codex round}

A v4 confirmation round is optional.  V4 only \emph{restricts} the
v3 claims (single-shock form unchanged; multi-shock scope narrowed
with explicit conservativeness argument), so a regression in
substantive content is unlikely.  The minor cleanups (Q1 single-shock
formalization, Q3 Poisson rate reframing, abstract fix) are
mechanical and verifiable by inspection.

If a v4 confirmation round runs, the questions to ask are:
\begin{enumerate}
\item Does the cooperative-overlap conservativeness argument
  actually hold? V4 argues equal-substrate analysis is a lower bound
  on $\Ddiv^{\mathrm{multi}}$ when overlap dominates redundancy; is
  this rigorous, or are there mixed regimes where the direction
  flips?
\item Is the $\kappa < 1 - \delta$ regime condition operationally
  meaningful? For canonical deployments with $\gamma \approx 1$ and
  realistic $\lambda$, what is the typical $\delta$ slack?
\item Does the cooperative-vs-redundancy audit have hidden
  ambiguities? Specifically, can a substrate-cooperative structure
  be classified differently under different threat models?
\end{enumerate}

Otherwise, Paper~10 may cite Lemma~5c v4 as established for
deployments meeting the cooperative-overlap regime condition.

\end{document}
